\chapter{2026年天津卷}

\section{单选题}

\begin{problem}
已知全集 \(U=\{-2,-1,0,1,2,3\}\)，集合 \(A=\{-1,0,1,3\}\)，集合 \(B=\{-2,0,1\}\)，则 \(\left( \complement_U A \right)\cup B=\)（\quad）

\choices
    {\(\{-2\}\)}
    {\(\{-2,2\}\)}
    {\(\{0,1,2\}\)}
    {\(\{-2,0,1,2\}\)}

\end{problem}

\begin{answer}
D
\end{answer}

\begin{solution}
因为 \(\complement_U A=U\setminus A=\{-2,2\}\)，所以 \(\left( \complement_U A \right)\cup B=\{-2,2\}\cup\{-2,0,1\}=\{-2,0,1,2\}\). 故选 \(D\).
\end{solution}
\begin{problem}
设 \(x\in \R\)，则“\(x>0\)”是“\(x^2+3x>0\)”的（\quad）

\choices
    {充分而不必要条件}
    {必要而不充分条件}
    {充分必要条件}
    {既不充分也不必要条件}

\end{problem}

\begin{answer}
A
\end{answer}

\begin{solution}
若 \(x>0\)，则 \(x+3>0\)，从而 \(x^2+3x=x\left( x+3 \right)>0\). 所以“\(x>0\)”是“\(x^2+3x>0\)”的充分条件.

由 \(x^2+3x>0\) 得 \(x\left( x+3 \right)>0\)，解得 \(x>0\) 或 \(x<-3\). 因此“\(x^2+3x>0\)”不能推出“\(x>0\)”，所以“\(x>0\)”不是必要条件.

故选 \(A\).
\end{solution}
\begin{problem}
调查候鸟和温度的关系，在不同温度下统计候鸟的数量，所得数据如图所示，其中相关系数 \(r=-0.91\)，根据最小二乘法算得：\(\hat{y}=-1.17x+1370.7\)，下列说法正确的是（\quad）

\bitmapfigure[max width=3.6cm,max totalheight=4.2cm]{img/2026/tianjin/q3_fig1.png}

\choices
    {\(y\) 与 \(x\) 负相关}
    {当 \(x=10\) 时，\(y\) 一定为 \(1359\)}
    {当 \(x=10\) 时，\(y\) 一定小于 \(1359\)}
    {两变量无线性关系}

\end{problem}

\begin{answer}
A
\end{answer}

\begin{solution}
由相关系数 \(r=-0.91\) 可知，\(r<0\) 且 \(|r|=0.91\) 较大，因此 \(x\) 与 \(y\) 具有较强的负线性相关关系，选项 \(A\) 正确，选项 \(D\) 错误.

当 \(x=10\) 时，\(\hat y=-1.17\times 10+1370.7=1359\). 这里的 \(1359\) 是回归方程给出的估计值，不是实际观测值，所以不能说 \(y\) 一定等于 \(1359\)，也不能说 \(y\) 一定小于 \(1359\). 因此选项 \(B\)，\(C\) 都错误.

故选 \(A\).
\end{solution}
\begin{problem}
函数 \(f(x)\) 的部分图象如图所示，则 \(f(x)\) 的解析式可能为（\quad）

\bitmapfigure[max width=4.2cm,max totalheight=4.6cm]{img/2026/tianjin/q4_fig1.png}

\choices
    {\(x+\sin \pi x\)}
    {\(x-\sin \pi x\)}
    {\(-x+\sin \pi x\)}
    {\(x\cdot \sin \pi x\)}

\end{problem}

\begin{answer}
C
\end{answer}

\begin{solution}
由图象可见，曲线过原点，且 \(x=1\) 时图象在 \(x\) 轴下方，所以 \(f(1)<0\).

四个选项在 \(x=1\) 处的函数值分别为：选项 \(A\) 为 \(1+\sin \pi=1\)，选项 \(B\) 为 \(1-\sin \pi=1\)，选项 \(C\) 为 \(-1+\sin \pi=-1\)，选项 \(D\) 为 \(1\cdot \sin \pi=0\). 只有选项 \(C\) 满足 \(f(1)<0\).

故选 \(C\).
\end{solution}
\begin{problem}
正方体 \(ABCD-A_1B_1C_1D_1\) 中，错误的是（\quad）

\choices
    {\(AC//A_1C_1\)}
    {\(CC_1\perp \text{面} ABC\)}
    {\(ACD_1//\text{面} A_1C_1B\)}
    {\(ADD_1\perp \text{面} ACD_1\)}

\end{problem}

\begin{answer}
D
\end{answer}

\begin{solution}
取正方体的棱长为 \(1\)，建立坐标系，可设 \(A(0,0,0)\)，\(B(1,0,0)\)，\(C(1,1,0)\)，\(D(0,1,0)\)，\(A_1(0,0,1)\)，\(B_1(1,0,1)\)，\(C_1(1,1,1)\)，\(D_1(0,1,1)\).

选项 \(A\) 中，\(\overrightarrow{AC}=\left( 1,1,0 \right)\)，\(\overrightarrow{A_1C_1}=\left( 1,1,0 \right)\)，所以 \(AC//A_1C_1\)，正确.

选项 \(B\) 中，\(CC_1\) 垂直于底面内两条相交直线 \(BC\) 与 \(CD\)，所以 \(CC_1\perp \text{面} ABC\)，正确.

选项 \(C\) 中，\(\overrightarrow{AC}=\left( 1,1,0 \right)\)，\(\overrightarrow{AD_1}=\left( 0,1,1 \right)\)，所以平面 \(ACD_1\) 的一个法向量可取为 \(\overrightarrow{AC}\times \overrightarrow{AD_1}=\left( 1,-1,1 \right)\). 又 \(\overrightarrow{A_1C_1}=\left( 1,1,0 \right)\)，\(\overrightarrow{A_1B}=\left( 1,0,-1 \right)\)，故平面 \(A_1C_1B\) 的一个法向量可取为 \(\overrightarrow{A_1C_1}\times \overrightarrow{A_1B}=\left( -1,1,-1 \right)\). 两个法向量平行，所以 \(\text{面} ACD_1//\text{面} A_1C_1B\)，正确.

选项 \(D\) 中，平面 \(ADD_1\) 是平面 \(x=0\)，其法向量可取为 \(\left( 1,0,0 \right)\). 由于 \(\left( 1,0,0 \right)\cdot \left( 1,-1,1 \right)=1\ne 0\)，所以 \(\text{面} ADD_1\) 与 \(\text{面} ACD_1\) 不垂直. 因而错误的是 \(D\).

故选 \(D\).
\end{solution}
\begin{problem}
已知函数 \(f(x)=|\ln x|\)，若 \(a=f\left( 2^{0.3} \right)\)，\(b=f\left( 3^{0.3} \right)\)，\(c=f\left( 3^{-0.5} \right)\)，则 \(a\)，\(b\)，\(c\) 的大小关系为（\quad）

\choices
    {\(a<b<c\)}
    {\(b<a<c\)}
    {\(c<b<a\)}
    {\(c<a<b\)}

\end{problem}

\begin{answer}
A
\end{answer}

\begin{solution}
因为 \(2^{0.3}>1\)，\(3^{0.3}>1\)，\(3^{-0.5}<1\)，所以 \(a=\left| \ln 2^{0.3} \right|=0.3\ln 2\)，\(b=\left| \ln 3^{0.3} \right|=0.3\ln 3\)，\(c=\left| \ln 3^{-0.5} \right|=0.5\ln 3\).

又因为 \(\ln 2<\ln 3\)，所以 \(0.3\ln 2<0.3\ln 3\)，即 \(a<b\).

并且 \(0.3\ln 3<0.5\ln 3\)，即 \(b<c\).

因此 \(a<b<c\).

故选 \(A\).
\end{solution}
\begin{problem}
\(\left( x+\frac{1}{x} \right)\left( x+\frac{4}{x} \right)\) 的最小值为（\quad）

\choices
    {\(10\)}
    {\(9\)}
    {\(8\)}
    {\(6\)}

\end{problem}

\begin{answer}
B
\end{answer}

\begin{solution}
因为 \(x\ne 0\)，所以 \(\left( x+\frac{1}{x} \right)\left( x+\frac{4}{x} \right)=x^2+5+\frac{4}{x^2}\). 又 \(x^2+\frac{4}{x^2}\ge 2\sqrt{x^2\cdot \frac{4}{x^2}}=4\)，因此 \(x^2+5+\frac{4}{x^2}\ge 9\).

当 \(x^2=2\) 时取等号，所以最小值为 \(9\).

故选 \(B\).
\end{solution}
\begin{problem}
已知 \(S_{2n}-S_n=n\)，\(a_3=6\)，则 \(a_7+a_8=\)（\quad）

\choices
    {\(68\)}
    {\(56\)}
    {\(-3\)}
    {\(-4\)}

\end{problem}

\begin{answer}
C
\end{answer}

\begin{solution}
由 \(S_{2n}-S_n=n\) 可知 \(a_{n+1}+a_{n+2}+\cdots +a_{2n}=n\).

令 \(n=2\)，得 \(a_3+a_4=2\). 又 \(a_3=6\)，所以 \(a_4=-4\).

令 \(n=3\)，得 \(a_4+a_5+a_6=3\).

令 \(n=4\)，得 \(a_5+a_6+a_7+a_8=4\).

后式减前式，得 \(a_7+a_8-a_4=1\). 代入 \(a_4=-4\)，得 \(a_7+a_8=-3\).

故选 \(C\).
\end{solution}
\begin{problem}
已知双曲线 \(\frac{x^2}{a^2}-\frac{y^2}{b^2}=1\)（\(a>0\)，\(b>0\)）的左焦点为 \(F\)，\(A\) 是右顶点，\(P\) 是双曲线上一点，满足 \(|FA|=|FP|\)，\(\angle FAP=30^\circ\)，则双曲线离心率为（\quad）

\choices
    {\(4\)}
    {\(\frac{8}{3}\)}
    {\(\frac{8}{5}\)}
    {\(\frac{4}{3}\)}

\end{problem}

\begin{answer}
D
\end{answer}

\begin{solution}
设双曲线的焦距为 \(2c\)，则左焦点 \(F\left( -c,0 \right)\)，右顶点 \(A\left( a,0 \right)\)，离心率 \(e=\frac{c}{a}\).

由 \(|FA|=|FP|\) 知，在 \(\triangle FAP\) 中，\(FA=FP\). 又 \(\angle FAP=30^\circ\)，所以 \(\angle FPA=30^\circ\)，从而 \(\angle AFP=120^\circ\). 由正弦定理得 \(\frac{AP}{\sin 120^\circ}=\frac{FA}{\sin 30^\circ}\)，即 \(AP=\sqrt{3}\,FA=\sqrt{3}\left( a+c \right)\).

不妨设 \(P\) 在上支上，则射线 \(AP\) 与 \(x\) 轴正向夹角为 \(150^\circ\). 因此 \(P\left( a-\frac{3}{2}\left( a+c \right),\frac{\sqrt{3}}{2}\left( a+c \right) \right)=\left( -\frac{a+3c}{2},\frac{\sqrt{3}\left( a+c \right)}{2} \right)\).

将 \(P\) 代入双曲线方程，得 \(\frac{\left( a+3c \right)^2}{4a^2}-\frac{3\left( a+c \right)^2}{4b^2}=1\). 又 \(c^2=a^2+b^2\)，所以 \(b^2=a^2\left( e^2-1 \right)\). 上式化为 \(\frac{\left( 1+3e \right)^2}{4}-\frac{3\left( 1+e \right)^2}{4\left( e^2-1 \right)}=1\). 整理得 \(3e^2+2e-1=\frac{e+1}{e-1}\)，于是 \(\left( 3e^2+2e-1 \right)\left( e-1 \right)=e+1\). 展开并整理，得 \(e\left( 3e^2-e-4 \right)=0\). 因为双曲线离心率 \(e>1\)，所以 \(3e^2-e-4=0\)，解得 \(e=\frac{4}{3}\).

故选 \(D\).
\end{solution}
\section{填空题}

\begin{problem}
化简 \((3+\i)^2=\fillinblank{}\).

\end{problem}

\begin{answer}
\(8+6\i\)
\end{answer}

\begin{solution}
由平方公式可得 \(\left( 3+\i \right)^2=9+6\i+\i^2=9+6\i-1=8+6\i\). 故答案为 \(8+6\i\).
\end{solution}
\begin{problem}
\((x+2y)^4\) 展开式中 \(x^3y\) 的系数为 \fillinblank{}.

\end{problem}

\begin{answer}
\(8\)
\end{answer}

\begin{solution}
由二项式定理，\(\left( x+2y \right)^4=\sum_{k=0}^4 \mathrm C_4^k x^{4-k}\left( 2y \right)^k\). 其中 \(x^3y\) 对应 \(k=1\)，故其系数为 \(\mathrm C_4^1\cdot 2=4\times 2=8\). 故答案为 \(8\).
\end{solution}
\begin{problem}
在 \(\triangle ABC\) 中，\(BC=4\)，\(AC=3\)，\(\cos A=-\frac{1}{4}\)，则 \(\sin B=\fillinblank{}\).

\end{problem}

\begin{answer}
\(\frac{3\sqrt{15}}{16}\)
\end{answer}

\begin{solution}
设 \(a=BC=4\)，\(b=AC=3\). 因为 \(\cos A=-\frac{1}{4}\)，所以 \(\sin A=\sqrt{1-\cos^2 A}=\sqrt{1-\frac{1}{16}}=\frac{\sqrt{15}}{4}\).

由正弦定理 \(\frac{a}{\sin A}=\frac{b}{\sin B}\)，得 \(\sin B=\frac{b\sin A}{a}=\frac{3\cdot \frac{\sqrt{15}}{4}}{4}=\frac{3\sqrt{15}}{16}\).

故答案为 \(\frac{3\sqrt{15}}{16}\).
\end{solution}
\begin{problem}
箱子里有一个红球，两个黄球，三个白球，有放回的取三次，三次都没取到黄球的概率是 \fillinblank{}；在三次都没取到黄球的条件下，至少取到一次红球的概率是 \fillinblank{}.

\end{problem}

\begin{answer}
\(\frac{8}{27}\)，\(\frac{37}{64}\)
\end{answer}

\begin{solution}
每次取球时，不取到黄球的概率为 \(\frac{1+3}{1+2+3}=\frac{4}{6}=\frac{2}{3}\). 因为是有放回抽取三次，所以三次都没取到黄球的概率为 \(\left( \frac{2}{3} \right)^3=\frac{8}{27}\).

在“三次都没取到黄球”的条件下，每次只可能取到红球或白球，此时单次取到红球的条件概率为 \(\frac{1}{4}\)，取到白球的条件概率为 \(\frac{3}{4}\). 因此至少取到一次红球的条件概率为 \(1-\left( \frac{3}{4} \right)^3=1-\frac{27}{64}=\frac{37}{64}\).

故答案为 \(\frac{8}{27}\)，\(\frac{37}{64}\).
\end{solution}
\begin{problem}
已知 \(|\bs{a}|=\bs{a}\cdot \bs{b}=1\)，\(|\bs{b}|>1\)，记 \(\bs{c}=\lambda \bs{a}+\mu \bs{b}\). 当 \(\bs{a}+\bs{b}-\bs{c}=0\) 时，\(\lambda+\mu=\fillinblank{}\)；当 \(|\bs{a}+\bs{b}-\bs{c}|=1\) 时，\(\lambda+\mu\) 的取值范围为 \fillinblank{}.

\end{problem}

\begin{answer}
\(2\)，\([1,3]\)
\end{answer}

\begin{solution}
由 \(\bs{c}=\lambda \bs{a}+\mu \bs{b}\) 得 \(\bs{a}+\bs{b}-\bs{c}=\left( 1-\lambda \right)\bs{a}+\left( 1-\mu \right)\bs{b}\).

当 \(\bs{a}+\bs{b}-\bs{c}=0\) 时，有 \(\left( 1-\lambda \right)\bs{a}+\left( 1-\mu \right)\bs{b}=0\). 由题意 \(\bs{a}\cdot \bs{b}=1\) 且 \(|\bs{a}|=1\)，又 \(|\bs{b}|>1\)，可知 \(\bs{a}\) 与 \(\bs{b}\) 不平行，所以 \(1-\lambda=0\)，\(1-\mu=0\). 于是 \(\lambda=\mu=1\)，故 \(\lambda+\mu=2\).

下面求第二空. 设 \(x=1-\lambda\)，\(y=1-\mu\)，则 \(\lambda+\mu=2-\left( x+y \right)\)，且 \(\left| \bs{a}+\bs{b}-\bs{c} \right|=1\) 等价于 \(\left| x\bs{a}+y\bs{b} \right|=1\). 两边平方，得 \(\left| x\bs{a}+y\bs{b} \right|^2=x^2|\bs{a}|^2+2xy\left( \bs{a}\cdot \bs{b} \right)+y^2|\bs{b}|^2=1\). 由 \(|\bs{a}|=1\)，\(\bs{a}\cdot \bs{b}=1\)，得 \(x^2+2xy+y^2|\bs{b}|^2=1\)，即 \(\left( x+y \right)^2+\left( |\bs{b}|^2-1 \right)y^2=1\). 因为 \(|\bs{b}|>1\)，所以 \(|\bs{b}|^2-1>0\)，从而 \(\left( x+y \right)^2\le 1\)，于是 \(-1\le x+y\le 1\). 所以 \(1\le \lambda+\mu=2-\left( x+y \right)\le 3\).

反过来，任取 \(t\in [1,3]\)，令 \(s=2-t\)，则 \(s\in [-1,1]\). 再取 \(y=\sqrt{\frac{1-s^2}{|\bs{b}|^2-1}}\)，\(x=s-y\)，则 \(\left( x+y \right)^2+\left( |\bs{b}|^2-1 \right)y^2=s^2+\left( 1-s^2 \right)=1\). 因此存在对应的 \(\lambda=1-x\)，\(\mu=1-y\)，使得 \(|\bs{a}+\bs{b}-\bs{c}|=1\) 且 \(\lambda+\mu=t\).

故 \(\lambda+\mu\) 的取值范围为 \([1,3]\).
\end{solution}
\begin{problem}
在平面内，\(O\) 为坐标原点，抛物线 \(y^2=2x\) 上有 \(A\)、\(B\)、\(C\)、\(D\) 四个点，\(A\)、\(B\)、\(C\)、\(D\) 的纵坐标分别为 \(y_A\)、\(y_B\)、\(y_C\)、\(y_D\)，直线 \(AB\) 与直线 \(CD\) 交 \(x\) 轴于点 \(P\)，直线 \(AC\) 交 \(x\) 轴于点 \(M\)，直线 \(BD\) 交 \(x\) 轴于点 \(N\)，以下说法正确的有 \fillinblank{}.

\begin{circlelist}
\item 若 \(P\) 与抛物线焦点重合，则 \(y_Ay_B=-2\)；
\item \(y_Ay_B=y_Cy_D\)；
\item \(|OM||ON|=2|OP|^2\)；
\item \(\left| y_A-y_C \right||OP|=\left| y_B-y_D \right||OM|\)；
\item \(\frac{S_{\triangle ACP}}{S_{\triangle BDP}}=\left( \frac{|OM|}{|ON|} \right)^2\)
\end{circlelist}

\end{problem}

\begin{answer}
\circled{2}\circled{4}
\end{answer}

\begin{solution}
设 \(y_A=u\)，\(y_B=v\)，\(y_C=w\)，\(y_D=z\)，则 \(A\left( \frac{u^2}{2},u \right)\)，\(B\left( \frac{v^2}{2},v \right)\)，\(C\left( \frac{w^2}{2},w \right)\)，\(D\left( \frac{z^2}{2},z \right)\).

对于抛物线 \(y^2=2x\) 上纵坐标分别为 \(m\)，\(n\) 的两点，连接所得弦与 \(x\) 轴的交点横坐标为 \(-\frac{mn}{2}\). 因此 \(P\left( -\frac{uv}{2},0 \right)=\left( -\frac{wz}{2},0 \right)\)，\(M\left( -\frac{uw}{2},0 \right)\)，\(N\left( -\frac{vz}{2},0 \right)\)，于是 \(uv=wz\).

先看 \circled{1}. 抛物线 \(y^2=2x\) 的焦点为 \(\left( \frac{1}{2},0 \right)\). 若 \(P\) 与焦点重合，则 \(-\frac{uv}{2}=\frac{1}{2}\)，所以 \(uv=-1\). 故 \(y_Ay_B=-1\)，不是 \(-2\)，所以 \circled{1} 错误.

由 \(uv=wz\) 立刻得到 \(y_Ay_B=y_Cy_D\)，所以 \circled{2} 正确.

又 \(|OP|=\frac{|uv|}{2}\)，\(|OM|=\frac{|uw|}{2}\)，\(|ON|=\frac{|vz|}{2}\). 因此 \(|OM||ON|=\frac{|uwvz|}{4}=\frac{|uv\cdot wz|}{4}=\frac{|uv|^2}{4}=|OP|^2\). 所以 \circled{3} 错误.

再看 \circled{4}. 由 \(uv=wz\) 得 \(z=\frac{uv}{w}\). 于是 \(\left| y_A-y_C \right||OP|=\left| u-w \right|\cdot \frac{|uv|}{2}\). 而 \(\left| y_B-y_D \right||OM|=\left| v-z \right|\cdot \frac{|uw|}{2}=\left| v-\frac{uv}{w} \right|\cdot \frac{|uw|}{2}=\left| \frac{v\left( w-u \right)}{w} \right|\cdot \frac{|uw|}{2}\). 化简得 \(\left| y_B-y_D \right||OM|=\left| u-w \right|\cdot \frac{|uv|}{2}\)，故 \(\left| y_A-y_C \right||OP|=\left| y_B-y_D \right||OM|\). 所以 \circled{4} 正确.

最后看 \circled{5}. 取一组满足条件的点，例如 \(u=1\)，\(v=2\)，\(w=-1\)，\(z=-2\). 此时 \(uv=wz=2\)，故两条直线 \(AB\) 与 \(CD\) 都交 \(x\) 轴于同一点 \(P\). 可得 \(A\left( \frac{1}{2},1 \right)\)，\(C\left( \frac{1}{2},-1 \right)\)，\(P(-1,0)\)，\(B(2,2)\)，\(D(2,-2)\). 于是 \(S_{\triangle ACP}=\frac{3}{2}\)，\(S_{\triangle BDP}=6\)，所以 \(\frac{S_{\triangle ACP}}{S_{\triangle BDP}}=\frac{1}{4}\). 又 \(|OM|=\frac{1}{2}\)，\(|ON|=2\)，从而 \(\left( \frac{|OM|}{|ON|} \right)^2=\left( \frac{1/2}{2} \right)^2=\frac{1}{16}\). 二者不相等，所以 \circled{5} 错误.

综上，正确的是 \(\circled{2}\circled{4}\).
\end{solution}
\section{解答题}

\begin{problem}
已知 \(f(x)=\sin \left( 2x+\frac{\pi}{6} \right)\).

\begin{enumerate}
\item 求最小正周期；
\item 若 \(x\in \left[ -\frac{\pi}{6},\frac{\pi}{12} \right]\)，求 \(f(x)\) 的最大值和最小值；
\item 若 \(\alpha\in \left( 0,\frac{\pi}{2} \right)\)，\(\sin \alpha=\frac{\sqrt{3}}{3}\)，求 \(f(\alpha)\).
\end{enumerate}

\end{problem}

\begin{solution}
（1）函数 \(f(x)=\sin \left( 2x+\frac{\pi}{6} \right)\) 的角频率为 \(2\)，所以最小正周期为 \(T=\frac{2\pi}{2}=\pi\).

（2）当 \(x\in \left[ -\frac{\pi}{6},\frac{\pi}{12} \right]\) 时，\(2x+\frac{\pi}{6}\in \left[ -\frac{\pi}{6},\frac{\pi}{3} \right]\). 因为正弦函数在区间 \(\left[ -\frac{\pi}{6},\frac{\pi}{3} \right]\) 上单调递增，所以 \(f_{\min}=\sin \left( -\frac{\pi}{6} \right)=-\frac{1}{2}\)，\(f_{\max}=\sin \left( \frac{\pi}{3} \right)=\frac{\sqrt{3}}{2}\).

（3）因为 \(\alpha\in \left( 0,\frac{\pi}{2} \right)\)，\(\sin \alpha=\frac{\sqrt{3}}{3}\)，所以 \(\cos \alpha=\sqrt{1-\sin^2 \alpha}=\sqrt{1-\frac{1}{3}}=\frac{\sqrt{6}}{3}\). 于是 \(\sin 2\alpha=2\sin \alpha \cos \alpha=2\cdot \frac{\sqrt{3}}{3}\cdot \frac{\sqrt{6}}{3}=\frac{2\sqrt{2}}{3}\)，\(\cos 2\alpha=\cos^2 \alpha-\sin^2 \alpha=\frac{2}{3}-\frac{1}{3}=\frac{1}{3}\). 从而 \(f(\alpha)=\sin \left( 2\alpha+\frac{\pi}{6} \right)=\sin 2\alpha \cos \frac{\pi}{6}+\cos 2\alpha \sin \frac{\pi}{6}\). 所以 \(f(\alpha)=\frac{2\sqrt{2}}{3}\cdot \frac{\sqrt{3}}{2}+\frac{1}{3}\cdot \frac{1}{2}=\frac{\sqrt{6}}{3}+\frac{1}{6}=\frac{2\sqrt{6}+1}{6}\).
\end{solution}
\begin{problem}
在长方体 \(ABCD-A_1B_1C_1D_1\) 中，\(AB=2\)，\(AA_1=3\)，\(AD=4\)，\(A_1E=3ED_1\)，\(2C_1F=FC\).

\begin{enumerate}
\item 求证：\(BD\perp \text{面} CEF\)；
\item 求面 \(AEF\) 与面 \(CEF\) 的夹角的余弦值；
\item 求三棱锥 \(A-CEF\) 的体积.
\end{enumerate}

\bitmapfigure[max width=4.5cm,max totalheight=4.8cm]{img/2026/tianjin/q17_fig1.png}

\end{problem}

\begin{solution}
（1）以 \(A\) 为原点，分别以 \(\overrightarrow{AB}\)，\(\overrightarrow{AD}\)，\(\overrightarrow{AA_1}\) 所在方向为 \(x\) 轴，\(y\) 轴，\(z\) 轴，建立空间直角坐标系，则 \(A(0,0,0)\)，\(B(2,0,0)\)，\(D(0,4,0)\)，\(C(2,4,0)\)，\(A_1(0,0,3)\)，\(C_1(2,4,3)\)，\(D_1(0,4,3)\).

由 \(A_1E=3ED_1\) 知 \(E\) 将 \(A_1D_1\) 按 \(3:1\) 分点，所以 \(E(0,3,3)\). 由 \(2C_1F=FC\) 知 \(C_1F:FC=1:2\)，所以 \(F(2,4,2)\).

于是 \(\overrightarrow{BD}=\left( -2,4,0 \right)\)，\(\overrightarrow{CE}=\left( -2,-1,3 \right)\)，\(\overrightarrow{CF}=\left( 0,0,2 \right)\). 计算得 \(\overrightarrow{CE}\times \overrightarrow{CF}=\left( -2,4,0 \right)\). 这说明 \(\overrightarrow{BD}\) 与平面 \(CEF\) 的一个法向量平行，所以 \(BD\perp \text{面} CEF\).

（2）面 \(CEF\) 的一个法向量可取为 \(\bs{n}_1=\overrightarrow{CE}\times \overrightarrow{CF}=\left( -2,4,0 \right)\).

又 \(\overrightarrow{AE}=\left( 0,3,3 \right)\)，\(\overrightarrow{AF}=\left( 2,4,2 \right)\)，所以面 \(AEF\) 的一个法向量可取为 \(\bs{n}_2=\overrightarrow{AE}\times \overrightarrow{AF}=\left( -6,6,-6 \right)\).

设两平面的夹角为 \(\theta\)，则 \(\cos \theta=\frac{\left| \bs{n}_1\cdot \bs{n}_2 \right|}{|\bs{n}_1|\cdot |\bs{n}_2|}\). 计算得 \(\bs{n}_1\cdot \bs{n}_2=\left( -2 \right)\left( -6 \right)+4\cdot 6=36\)，且 \(|\bs{n}_1|=\sqrt{20}=2\sqrt{5}\)，\(|\bs{n}_2|=\sqrt{108}=6\sqrt{3}\). 因此 \(\cos \theta=\frac{36}{2\sqrt{5}\cdot 6\sqrt{3}}=\frac{\sqrt{15}}{5}\).

（3）三棱锥 \(A-CEF\) 的体积为 \(V=\frac{1}{6}\left| \det \left( \overrightarrow{AC},\overrightarrow{AE},\overrightarrow{AF} \right) \right|\). 其中 \(\overrightarrow{AC}=\left( 2,4,0 \right)\)，\(\overrightarrow{AE}=\left( 0,3,3 \right)\)，\(\overrightarrow{AF}=\left( 2,4,2 \right)\). 于是
\[
\det \left( \overrightarrow{AC},\overrightarrow{AE},\overrightarrow{AF} \right)
=
\begin{vmatrix}
2&4&0\\
0&3&3\\
2&4&2
\end{vmatrix}
=12
\]
所以 \(V=\frac{12}{6}=2\).
\end{solution}
\begin{problem}
已知椭圆 \(C:\frac{x^2}{a^2}+\frac{y^2}{b^2}=1\)（\(a>b>0\)）的离心率为 \(\frac{1}{2}\)，椭圆被直线 \(x=b\) 截得的线段长为 \(\sqrt{3}\).

\begin{enumerate}
\item 求 \(C\) 的标准方程；
\item 斜率为 \(-\sqrt{3}\) 的直线与圆 \(x^2+y^2=b^2\) 相切，且该直线交椭圆于 \(P(x_1,y_1)\)，\(Q(x_2,y_2)\)（\(y_1<y_2\)），\(A\) 是椭圆的上顶点. 记直线 \(AP\)，\(AQ\) 的斜率分别为 \(k_1\)，\(k_2\)，求 \(\frac{k_1}{k_2}\).
\end{enumerate}

\end{problem}

\begin{solution}
（1）设椭圆的焦距为 \(2c\). 由离心率为 \(\frac{1}{2}\)，得 \(\frac{c}{a}=\frac{1}{2}\)，所以 \(c^2=\frac{a^2}{4}\). 又椭圆满足 \(c^2=a^2-b^2\)，故 \(a^2-b^2=\frac{a^2}{4}\)，即 \(b^2=\frac{3a^2}{4}\).

因为直线 \(x=b\) 与椭圆交于两点，把 \(x=b\) 代入椭圆方程，得 \(\frac{b^2}{a^2}+\frac{y^2}{b^2}=1\). 由 \(\frac{b^2}{a^2}=\frac{3}{4}\)，得 \(\frac{y^2}{b^2}=\frac{1}{4}\)，所以 \(y=\pm \frac{b}{2}\). 该弦长为 \(\frac{b}{2}-\left( -\frac{b}{2} \right)=b\). 已知弦长为 \(\sqrt{3}\)，故 \(b=\sqrt{3}\)，即 \(b^2=3\). 于是 \(a^2=\frac{4}{3}b^2=4\).

因此椭圆 \(C\) 的标准方程为 \(\frac{x^2}{4}+\frac{y^2}{3}=1\).

（2）由第（1）问知 \(b^2=3\)，所给圆为 \(x^2+y^2=3\). 设所求切线为 \(y=-\sqrt{3}x+t\). 因为该直线与圆相切，所以原点到直线的距离等于半径 \(\sqrt{3}\)，即 \(\frac{|t|}{\sqrt{1+3}}=\sqrt{3}\)，解得 \(|t|=2\sqrt{3}\). 所以切线只有两条：\(y=-\sqrt{3}x+2\sqrt{3}\) 或 \(y=-\sqrt{3}x-2\sqrt{3}\).

先取 \(y=-\sqrt{3}x+2\sqrt{3}\). 代入椭圆 \(\frac{x^2}{4}+\frac{y^2}{3}=1\)，得 \(\frac{x^2}{4}+\frac{\left( -\sqrt{3}x+2\sqrt{3} \right)^2}{3}=1\). 化简得 \(5x^2-16x+12=0\)，解得 \(x=2\) 或 \(x=\frac{6}{5}\). 对应的两点为 \(P(2,0)\)，\(Q\left( \frac{6}{5},\frac{4\sqrt{3}}{5} \right)\)，且满足 \(y_1<y_2\). 又上顶点 \(A(0,\sqrt{3})\). 于是 \(k_1=\frac{0-\sqrt{3}}{2-0}=-\frac{\sqrt{3}}{2}\)，\(k_2=\frac{\frac{4\sqrt{3}}{5}-\sqrt{3}}{\frac{6}{5}-0}=-\frac{\sqrt{3}}{6}\)，所以 \(\frac{k_1}{k_2}=3\).

若取另一条切线 \(y=-\sqrt{3}x-2\sqrt{3}\)，把上面所求两交点的横、纵坐标同时取相反数，可得交点为 \(\left( -\frac{6}{5},-\frac{4\sqrt{3}}{5} \right)\) 和 \((-2,0)\)，对应两条直线 \(AP\)，\(AQ\) 的斜率之比仍为 \(3\).

故 \(\frac{k_1}{k_2}=3\).
\end{solution}
\begin{problem}
已知等差数列 \(\{a_n\}\) 与等比数列 \(\{b_n\}\) 满足：\(a_1=2\)，\(b_1=1\)，\(a_2=b_1+b_2\)，\(a_4=b_3-b_1\).

\begin{enumerate}
\item 求数列 \(\{a_n\}\)，\(\{b_n\}\) 的通项公式；
\item 记 \(E_n=\left\{ x\in \R\mid x\le n,\ \exists k\in \N^*,\ x=a_k\text{ 或 }x=b_k \right\}\)，记 \(c_n\) 为 \(E_n\) 中的元素个数.

    \begin{enumerate}
    \item 求 \(c_{3^n}\)；
    \item 求 \(\sum_{m=1}^{3^n-1}(-1)^m\cdot a_m\cdot c_m\).
    \end{enumerate}
\end{enumerate}

\end{problem}

\begin{solution}
（1）设等差数列 \(\{a_n\}\) 的公差为 \(d\)，等比数列 \(\{b_n\}\) 的公比为 \(q\). 则 \(a_n=2+\left( n-1 \right)d\)，\(b_n=q^{n-1}\).

由 \(a_2=b_1+b_2\)，得 \(2+d=1+q\)，所以 \(d=q-1\). 由 \(a_4=b_3-b_1\)，得 \(2+3d=q^2-1\). 代入 \(d=q-1\)，得 \(2+3\left( q-1 \right)=q^2-1\)，即 \(q^2-3q=0\). 因此 \(q=0\) 或 \(q=3\).

若 \(q=0\)，则 \(b_1=1\)，\(b_2=0\). 这时从第二项起数列中出现零项，不能满足高中阶段等比数列“相邻两项之比为同一个非零常数”的通常定义，所以舍去.

故 \(q=3\)，从而 \(d=2\). 于是 \(a_n=2+\left( n-1 \right)\cdot 2=2n\)，\(b_n=3^{n-1}\).

（2）（i）由第（1）问知 \(a_k=2k\)，\(b_k=3^{k-1}\).

因为 \(a_k\le n\) 等价于 \(k\le \frac{n}{2}\)，所以满足条件的 \(a_k\) 共 \(\left\lfloor \frac{n}{2} \right\rfloor\) 个.

又因为 \(b_k\le n\) 等价于 \(3^{k-1}\le n\)，即 \(k\le \left\lfloor \log_3 n \right\rfloor +1\)，所以满足条件的 \(b_k\) 共 \(\left\lfloor \log_3 n \right\rfloor +1\) 个.

并且 \(a_k=2k\) 全是正偶数，\(b_k=3^{k-1}\) 全是正奇数，所以两部分没有重复元素. 故 \(c_n=\left\lfloor \frac{n}{2} \right\rfloor +\left\lfloor \log_3 n \right\rfloor +1\).

令 \(n\) 替换为 \(3^n\)，因为 \(3^n\) 为奇数，所以 \(\left\lfloor \frac{3^n}{2} \right\rfloor =\frac{3^n-1}{2}\)，且 \(\left\lfloor \log_3 3^n \right\rfloor =n\). 因此 \(c_{3^n}=\frac{3^n-1}{2}+n+1=\frac{3^n+2n+1}{2}\).

（ii）由上式，\(c_m=\left\lfloor \frac{m}{2} \right\rfloor +\left\lfloor \log_3 m \right\rfloor +1\). 又 \(a_m=2m\)，所以
\[
\sum_{m=1}^{3^n-1}\left( -1 \right)^m a_m c_m
=2\sum_{m=1}^{3^n-1}\left( -1 \right)^m m\left( \left\lfloor \frac{m}{2} \right\rfloor +\left\lfloor \log_3 m \right\rfloor +1 \right)
\]
记 \(N=3^n-1\)，\(K=\frac{N}{2}=\frac{3^n-1}{2}\). 将上式拆成三部分，记 \(S=2\left( T_1+T_2+T_3 \right)\)，其中 \(T_1=\sum_{m=1}^{N}\left( -1 \right)^m m\left\lfloor \frac{m}{2} \right\rfloor\)，\(T_2=\sum_{m=1}^{N}\left( -1 \right)^m m\)，\(T_3=\sum_{m=1}^{N}\left( -1 \right)^m m\left\lfloor \log_3 m \right\rfloor\).

先求 \(T_1\). 按 \(m=2k-1\) 与 \(m=2k\) 配对，有
\[
T_1=\sum_{k=1}^{K}\left[ -\left( 2k-1 \right)\left( k-1 \right)+2k^2 \right]
=\sum_{k=1}^{K}\left( 3k-1 \right)
\]
于是 \(T_1=\frac{3K\left( K+1 \right)}{2}-K=\frac{\left( 3^n-1 \right)\left( 3^{n+1}-1 \right)}{8}\).

再求 \(T_2\). 仍按奇偶配对，\(T_2=\sum_{k=1}^{K}\left[ -\left( 2k-1 \right)+2k \right]=\sum_{k=1}^{K}1=K=\frac{3^n-1}{2}\).

最后求 \(T_3\). 对于 \(j=0,1,\cdots,n-1\)，当 \(3^j\le m\le 3^{j+1}-1\) 时，有 \(\left\lfloor \log_3 m \right\rfloor =j\). 所以 \(T_3=\sum_{j=0}^{n-1}j\sum_{m=3^j}^{3^{j+1}-1}\left( -1 \right)^m m\). 注意到 \(\sum_{m=1}^{2r}\left( -1 \right)^m m=r\)，而 \(3^j-1\) 与 \(3^{j+1}-1\) 都是偶数，因此 \(\sum_{m=3^j}^{3^{j+1}-1}\left( -1 \right)^m m=\frac{3^{j+1}-1}{2}-\frac{3^j-1}{2}=3^j\). 故 \(T_3=\sum_{j=0}^{n-1}j\cdot 3^j\)，这是等比型求和，结果为 \(T_3=\frac{3-n3^n+\left( n-1 \right)3^{n+1}}{4}\).

于是
\[
T_1+T_2=\frac{\left( 3^n-1 \right)\left( 3^{n+1}-1 \right)}{8}+\frac{3^n-1}{2}
=\frac{3\left( 3^{2n}-1 \right)}{8}
\]
所以
\[
S=2\left[ \frac{3\left( 3^{2n}-1 \right)}{8}+\frac{3-n3^n+\left( n-1 \right)3^{n+1}}{4} \right]
\]
化简得 \(S=\frac{3^{2n+1}+\left( 4n-6 \right)3^n+3}{4}\).

故
\[
\sum_{m=1}^{3^n-1}\left( -1 \right)^m a_m c_m=\frac{3^{2n+1}+\left( 4n-6 \right)3^n+3}{4}
\]
\end{solution}
\begin{problem}
已知 \(f(x)=\e^x-\frac{2}{3}\sin x\).

\begin{enumerate}
\item 求曲线 \(y=f(x)\) 在点 \((0,f(0))\) 处的切线方程；
\item 当 \(x\in \left[ -\frac{1}{3},+\infty \right)\) 时，证明 \(f(x)\ge 1+\frac{x}{3}\)；
\item 求实数 \(a\) 的最大可能值，使得 \(f\left( \frac{1}{1} \right)\cdot f\left( \frac{1}{2} \right)\cdot f\left( \frac{1}{3} \right)\cdots f\left( \frac{1}{n} \right)\ge (n+1)^a\) 对任意的 \(n\in \N^*\) 都成立.
\end{enumerate}

\end{problem}

\begin{solution}
（1）因为 \(f(0)=\e^0-\frac{2}{3}\sin 0=1\)，且 \(f'(x)=\e^x-\frac{2}{3}\cos x\)，所以 \(f'(0)=1-\frac{2}{3}=\frac{1}{3}\). 故所求切线方程为 \(y-1=\frac{1}{3}\left( x-0 \right)\)，即 \(y=1+\frac{x}{3}\).

（2）设 \(g(x)=f(x)-1-\frac{x}{3}=\e^x-\frac{2}{3}\sin x-1-\frac{x}{3}\). 则 \(g'(x)=\e^x-\frac{2}{3}\cos x-\frac{1}{3}\)，\(g''(x)=\e^x+\frac{2}{3}\sin x\).

当 \(x\in \left[ -\frac{1}{3},0 \right]\) 时，由 \(\sin x\ge -\frac{1}{3}\) 得 \(g''(x)\ge \e^{-1/3}-\frac{2}{9}>0\). 当 \(x\ge 0\) 时，由 \(\sin x\ge -1\) 得 \(g''(x)\ge \e^x-\frac{2}{3}\ge \frac{1}{3}>0\). 所以当 \(x\ge -\frac{1}{3}\) 时，\(g''(x)>0\)，从而 \(g'(x)\) 严格递增.

又 \(g'(0)=1-\frac{2}{3}-\frac{1}{3}=0\)，故当 \(x\in \left[ -\frac{1}{3},0 \right]\) 时，\(g'(x)\le 0\)；当 \(x\ge 0\) 时，\(g'(x)\ge 0\). 于是 \(g(x)\) 在 \(x=0\) 处取得最小值，而 \(g(0)=0\)，所以 \(g(x)\ge 0\). 即当 \(x\ge -\frac{1}{3}\) 时，
\[
f(x)\ge 1+\frac{x}{3}
\]

（3）记 \(P_n=f\left( \frac{1}{1} \right)f\left( \frac{1}{2} \right)\cdots f\left( \frac{1}{n} \right)\).

由第（2）问可知，对任意 \(k\in \N^*\)，都有 \(f\left( \frac{1}{k} \right)\ge 1+\frac{1}{3k}\). 因此 \(P_n\ge \prod_{k=1}^{n}\left( 1+\frac{1}{3k} \right)\).

又因为 \(\left( 1+\frac{1}{3k} \right)^3=1+\frac{1}{k}+\frac{1}{3k^2}+\frac{1}{27k^3}\ge 1+\frac{1}{k}\)，所以
\[
\prod_{k=1}^{n}\left( 1+\frac{1}{3k} \right)^3\ge \prod_{k=1}^{n}\left( 1+\frac{1}{k} \right)=n+1
\]
从而 \(P_n^3\ge n+1\)，即 \(P_n\ge \left( n+1 \right)^{1/3}\). 故 \(a=\frac{1}{3}\) 可以取到.

下面证明 \(a\) 不可能大于 \(\frac{1}{3}\).

先证明当 \(0<x\le 1\) 时，存在常数 \(B=\e+\frac{1}{3}\)，使 \(f(x)\le 1+\frac{x}{3}+Bx^2\).

设 \(u(x)=1+x+\e x^2-\e^x\). 当 \(0\le x\le 1\) 时，\(u''(x)=2\e-\e^x\ge \e>0\). 又 \(u'(0)=0\)，\(u(0)=0\)，所以 \(u(x)\ge 0\)，即 \(\e^x\le 1+x+\e x^2\).

设 \(v(x)=\sin x-x+\frac{x^2}{2}\). 当 \(0\le x\le 1\) 时，\(v''(x)=1-\sin x\ge 0\). 又 \(v'(0)=0\)，\(v(0)=0\)，所以 \(v(x)\ge 0\)，即 \(\sin x\ge x-\frac{x^2}{2}\). 因此 \(-\frac{2}{3}\sin x\le -\frac{2}{3}x+\frac{x^2}{3}\). 于是 \(f(x)\le 1+\frac{x}{3}+\left( \e+\frac{1}{3} \right)x^2\).

令 \(x=\frac{1}{k}\)，则对任意 \(k\in \N^*\)，有 \(f\left( \frac{1}{k} \right)\le 1+\frac{1}{3k}+\frac{B}{k^2}\). 因此
\[
\ln P_n=\sum_{k=1}^{n}\ln f\left( \frac{1}{k} \right)
\le \sum_{k=1}^{n}\ln \left( 1+\frac{1}{3k}+\frac{B}{k^2} \right)
\]
由 \(\ln \left( 1+t \right)\le t\)（\(t>-1\)），得 \(\ln P_n\le \frac{1}{3}\sum_{k=1}^{n}\frac{1}{k}+B\sum_{k=1}^{n}\frac{1}{k^2}\).

又 \(\sum_{k=1}^{n}\frac{1}{k}\le 1+\ln n\le 1+\ln \left( n+1 \right)\). 对 \(k\ge 2\)，有 \(\frac{1}{k^2}\le \frac{1}{k\left( k-1 \right)}=\frac{1}{k-1}-\frac{1}{k}\)，所以对任意 \(n\in \N^*\)，都有 \(\sum_{k=1}^{n}\frac{1}{k^2}\le 2\). 因此
\[
\ln P_n\le \frac{1}{3}\ln \left( n+1 \right)+\frac{1}{3}+2B
\]
记 \(C=\e^{\frac{1}{3}+2B}\)，则 \(P_n\le C\left( n+1 \right)^{1/3}\).

若 \(a>\frac{1}{3}\)，取正整数 \(n\) 满足 \(n+1>C^{\frac{1}{a-\frac{1}{3}}}\). 这时 \(\left( n+1 \right)^a>C\left( n+1 \right)^{1/3}\ge P_n\)，与题意矛盾.

因此实数 \(a\) 的最大可能值为 \(\frac{1}{3}\).
\end{solution}
